\subsection{Algoritmo de Mo}

\begin{frame}
  \frametitle{Consultas de caminos usando el truco de Mo}

  Esta idea es muy buena pero depende de propiedades particulares de la operación \texttt{xor}:

  \vspace{0.5em}

  \begin{itemize}
    \item es asociativa
    \item tiene elemento neutro
    \item tiene inverso para cada elemento
    \item es conmutativa
    \item \textbf{cada elemento es su propio inverso}
  \end{itemize}
\end{frame}

\begin{frame}
  \frametitle{Consultas de caminos usando el truco de Mo: generalización}

  En realidad se puede generalizar a cualquier operación si cumple las primeras 4.

  \vspace{0.75em}

  \begin{alertblock}{Nota}
    conmutativa es opcional pero complica la implementación
  \end{alertblock}
\end{frame}

\begin{frame}
  \frametitle{Consultas de caminos usando el truco de Mo: idea}

  Básicamente la idea es hacer un \texttt{for} desde la posición de \(u\) hasta la de \(v\), e ir acumulando los valores.

  \vspace{0.75em}

  A medida que hacemos esto, hay que chequear si estamos yendo o volviendo por una arista. En caso de estar yendo, acumulamos el valor de la arista. En caso de estar volviendo, acumulamos el inverso del valor de la arista.

  \vspace{0.75em}

  De esta forma, tenemos un algoritmo que funciona para una operación donde no necesariamente cada elemento es su propio inverso.
\end{frame}

\begin{frame}
  \frametitle{Consultas de caminos usando el truco de Mo: implementación}

  Para chequear esto, simplemente guardamos los índices de las aristas en el Euler tour en vez de los valores.

  \vspace{0.75em}

  Aparte, podemos mantener un array de booleanos que nos diga si ya sumamos el valor de la arista o no.
\end{frame}

\begin{frame}[fragile]
  \frametitle{Consultas de caminos usando el truco de Mo: código}

\begin{lstlisting}
bool tengo[maxn];
int query(int u, int v) {
    int ans = 0;
    int l = position[u], r = position[v];
    if (l > r) swap(l, r);
    for (int i = l; i < r; i++) {
        int j = euler_tour[i];
        if (tengo[j]) ans -= edges[j].x;
        else          ans += edges[j].x;
        tengo[j] = !tengo[j];
    }
    return ans;
}
\end{lstlisting}
\end{frame}

\begin{frame}
  \frametitle{Consultas de caminos usando el truco de Mo: complejidad}

  Esto funciona pero es horrendamente lento. Claramente hace un \texttt{for} \(O(N)\) por consulta.
\end{frame}
